\documentclass[11pt]{article}
\usepackage[margin=2.5cm]{geometry}
\usepackage{lmodern}
\usepackage[T1]{fontenc}
\usepackage{microtype}
\usepackage{amsmath,amssymb}
\usepackage{booktabs}
\usepackage{algorithm}
\usepackage{algpseudocode}
\usepackage[hidelinks]{hyperref}

\newcommand{\slh}{SLH-DSA}
\newcommand{\lean}{leanSPHINCS}
\newcommand{\wc}{WOTS$^+$C}
\newcommand{\bytes}[1]{#1\,\text{B}}
\newcommand{\concat}{\,\|\,}

\title{\lean{}: an \slh{} variant for $2^{30}$ signatures\\with a cheap pruned signing mode}
\author{}
\date{Draft --- 30 September 2026}

\begin{document}
\maketitle

\section{Overview}

\lean{} is a stateless hash-based signature scheme at NIST security category~1 ($n = 16$ bytes).
It keeps the shape of \slh{}-SHA2-128-24 from NIST SP~800-230~\cite{sp800-230}: a single XMSS tree
($d = 1$) whose leaves are one-time keys, each certifying a FORS few-time key that signs the message.
It changes three things.
\begin{enumerate}
  \item \textbf{Parameters.} $h = 26$, $a = 10$, $k = 24$ (instead of $22, 24, 6$). A key can sign
        $2^{30}$ messages at 128-bit security, instead of $2^{24}$.
  \item \textbf{\wc{} instead of WOTS$^+$}~\cite{sphincsc}. The one-time keys drop the checksum chains.
        The signer instead grinds a 32-bit counter until the message digits have a fixed sum. This saves
        \bytes{60} and makes verification cost constant.
  \item \textbf{Pruned keys (signer-side only).} A weak signer may generate only a few thousand of the
        $2^{26}$ leaves, in about $10^6$ compressions, and still sign $2^{16}$ messages. Their signatures are
        ordinary \lean{} signatures; the verifier cannot tell the difference.
\end{enumerate}
Everything else --- FORS, the XMSS tree, the hash functions, addresses and the message hash --- is FIPS~205
\slh{}~\cite{fips205} unchanged.

\begin{table}[ht]
\centering
\begin{tabular}{cccccccccccc}
\toprule
$n$ & $h$ & $d$ & $a$ & $k$ & $w$ & chains & $S$ & $m$ & pk & sk & signature \\
\midrule
16 & 26 & 1 & 10 & 24 & 4 & 64 & 120 & 34 & \bytes{32} & \bytes{64} & \bytes{5\,684} \\
\bottomrule
\end{tabular}
\caption{Parameters. $w$, chains and $S$ describe \wc{}; $m$ is the message-digest length in bytes.}
\label{tab:params}
\end{table}

\begin{table}[ht]
\centering
\begin{tabular}{lr}
\toprule
Lifetime at 128-bit security & $2^{30.11}$ signatures \\
Verification (every signature) & 366 hash calls, 392 compressions \\
Full key: key generation / signing & $18.4\cdot10^{9}$ / $4.2\cdot10^{5}$ compressions \\
Pruned key ($2^{16}$ signatures): key generation / signing & $1.04\cdot10^{6}$ / $2.0\cdot10^{5}$ compressions \\
Signer cache (built at key generation) & 1~MiB \\
\bottomrule
\end{tabular}
\caption{Headline figures. Costs are SHA-256 compression calls (Section~\ref{sec:costs}).}
\label{tab:headline}
\end{table}

\section{Specification}

The scheme uses the FIPS~205 algorithms and the SHA2 security-category-1 instantiation of $\mathbf{F}$,
$\mathbf{H}$, $\mathbf{T}_\ell$, $\mathbf{PRF}$, $\mathbf{PRF}_{msg}$ and $\mathbf{H}_{msg}$
(FIPS~205 \S11.2.1). A SHAKE variant is obtained the same way from \S11.1. This section lists only what
differs.

\paragraph{Keys.} As in FIPS~205:
$\mathsf{pk} = (\mathsf{PK.seed}, \mathsf{PK.root})$ and
$\mathsf{sk} = (\mathsf{SK.seed}, \mathsf{SK.prf}, \mathsf{PK.seed}, \mathsf{PK.root})$.
$\mathsf{PK.root}$ is the root of one XMSS tree of height $h = 26$ whose leaves are \wc{} public keys.

\paragraph{Message digest.} $\mathit{digest} = \mathbf{H}_{msg}(R, \mathsf{PK.seed}, \mathsf{PK.root}, M')$
has $m = 34$ bytes. The first $\lceil ka/8 \rceil = 30$ bytes give the 24 FORS indices of 10 bits. The next
4 bytes, reduced mod $2^{26}$, give the leaf index $\mathit{idx}$. There is no tree index, since $d = 1$.

\paragraph{Signature.}
\[
  \mathsf{SIG} = R \concat \mathsf{SIG}_{FORS} \concat \mathit{ctr} \concat \mathsf{SIG}_{WOTS} \concat \mathsf{AUTH}
\]
with sizes $16 + 24\cdot11\cdot16 + 4 + 64\cdot16 + 26\cdot16 = \bytes{5\,684}$. Here $\mathsf{SIG}_{FORS}$ is
the FORS signature ($k = 24$ secret values and $a = 10$ authentication nodes each), $\mathit{ctr}$ is the
32-bit \wc{} counter, $\mathsf{SIG}_{WOTS}$ holds the 64 \wc{} chain values and $\mathsf{AUTH}$ is the XMSS
authentication path of leaf $\mathit{idx}$.

\subsection{\texorpdfstring{\wc{}}{WOTS+C}}

\wc{} uses $\ell = 64$ hash chains of length $w = 4$ ($\lg w = 2$) and no checksum chains. The message is
the $n$-byte FORS public key, so it is always the same for a given leaf. The counter search selects a
digest whose 64 base-4 digits sum to the fixed value $S = 120$:
\[
  \mathit{dig}(M, \mathit{ctr}) = \mathit{base\_2b}\bigl(\mathbf{T}_C(\mathsf{PK.seed}, \mathsf{ADRS}_C,
  M \concat \mathrm{toByte}(\mathit{ctr}, 4)),\, 2,\, 64\bigr),
\]
where, for SHA2,
\[
  \mathbf{T}_C(\mathsf{PK.seed}, \mathsf{ADRS}, X) = \mathrm{Trunc}_n\bigl(\mathrm{SHA\text{-}256}
  (\mathsf{PK.seed} \concat \mathrm{toByte}(0, 64 - n) \concat \mathsf{ADRS}^c \concat X)\bigr).
\]
$\mathsf{ADRS}_C$ is the leaf's WOTS address with a new type value \textsc{wots\_c} $= 7$.

\begin{algorithm}[ht]
\caption{\wc{} signing and public-key recovery (chain, PRF and $\mathbf{T}_{64}$ as in FIPS~205)}
\begin{algorithmic}[1]
\Function{wotsc\_sign}{$M$, $\mathsf{SK.seed}$, $\mathsf{PK.seed}$, $\mathsf{ADRS}$}
  \State $\mathit{ctr} \gets 0$
  \While{$\sum_i \mathit{dig}(M, \mathit{ctr})_i \neq S$} \Comment{about 829 tries for $S = 120$}
    \State $\mathit{ctr} \gets \mathit{ctr} + 1$
  \EndWhile
  \State $d \gets \mathit{dig}(M, \mathit{ctr})$
  \For{$i = 0, \dots, 63$}
    \State $\sigma_i \gets \mathrm{chain}(\mathbf{PRF}(\mathsf{PK.seed}, \mathsf{SK.seed}, \mathsf{ADRS}_i), 0, d_i)$
  \EndFor
  \State \Return $(\mathit{ctr}, \sigma_0, \dots, \sigma_{63})$
\EndFunction
\Statex
\Function{wotsc\_pkFromSig}{$(\mathit{ctr}, \sigma)$, $M$, $\mathsf{PK.seed}$, $\mathsf{ADRS}$}
  \State $d \gets \mathit{dig}(M, \mathit{ctr})$
  \If{$\sum_i d_i \neq S$} \Return $\bot$ \Comment{reject the signature}
  \EndIf
  \For{$i = 0, \dots, 63$}
    \State $p_i \gets \mathrm{chain}(\sigma_i, d_i, 3 - d_i)$
  \EndFor
  \State \Return $\mathbf{T}_{64}(\mathsf{PK.seed}, \mathsf{ADRS}_{pk}, p_0 \concat \dots \concat p_{63})$
\EndFunction
\end{algorithmic}
\end{algorithm}

The signer walks $\sum_i d_i = S = 120$ chain steps and the verifier walks $\sum_i (3 - d_i) = 192 - S = 72$,
for every signature. Table~\ref{tab:S} shows the trade-off behind $S = 120$: the counter search is
negligible, and a larger $S$ would save few verifier hashes for much more search.

\begin{table}[ht]
\centering
\begin{tabular}{rrrr}
\toprule
$S$ & valid digests & expected counter tries & verifier chain steps \\
\midrule
96  & $2^{123.5}$ & 22 & 96 \\
112 & $2^{121.2}$ & 110 & 80 \\
\textbf{120} & $\mathbf{2^{118.3}}$ & \textbf{829} & \textbf{72} \\
128 & $2^{114.2}$ & 14\,670 & 64 \\
132 & $2^{111.6}$ & 86\,554 & 60 \\
\bottomrule
\end{tabular}
\caption{Choice of the \wc{} digit sum $S$ (64 digits in base 4).}
\label{tab:S}
\end{table}

\subsection{Verification}

Verification is FIPS~205 \slh{} verification with $d = 1$, except that the one-time signature is checked with
\textsc{wotsc\_pkFromSig}, which rejects if the digit sum is not $S$. Its cost does not depend on the
message:
\begin{center}
\begin{tabular}{lrr}
\toprule
Step & hash calls & compressions \\
\midrule
$\mathbf{H}_{msg}$ & 1 & 5 \\
FORS: 24 trees $\times$ (1 leaf + 10 nodes) & 264 & 264 \\
FORS roots, $\mathbf{T}_{24}$ & 1 & 7 \\
\wc{} digest and chains & $1 + 72$ & $1 + 72$ \\
\wc{} public key, $\mathbf{T}_{64}$ & 1 & 17 \\
XMSS path & 26 & 26 \\
\midrule
Total & 366 & 392 \\
\bottomrule
\end{tabular}
\end{center}

\section{Security}

\paragraph{Few-time layer.} The message digest picks a leaf, and so a FORS instance, uniformly among
$2^{h}$. After $q$ signatures an instance has been used $r$ times with probability close to
$\mathrm{Pois}(\lambda, r)$, where $\lambda = q/2^h$. A forger wins on a random digest if all $k$ of its FORS indices
point to leaves already revealed. The generic-attack estimate used for SPHINCS$^+$, \slh{} and
SP~800-230~\cite{sphincsplus,fenner,kudinov-nick} is
\[
  \Pr[\text{forgery per query}] \;=\; \sum_{r \ge 1} \mathrm{Pois}(\lambda, r)\,
  \bigl(1 - (1 - 2^{-a})^{r}\bigr)^{k},
\]
and the security level is $\min(128, -\log_2 \Pr)$. With $h = 26$, $a = 10$ and $k = 24$, the scheme keeps
128 bits up to $2^{30.11}$ signatures:
\begin{center}
\begin{tabular}{lcccccc}
\toprule
Signatures & $2^{28}$ & $2^{29}$ & $2^{30}$ & $2^{31}$ & $2^{32}$ & $2^{34}$ \\
Security (bits) & 128 & 128 & 128 & 111.9 & 92.3 & 51.1 \\
\bottomrule
\end{tabular}
\end{center}
Security falls quickly past the limit, by about 18 bits per doubling, because each signature reveals
24 secrets. Signers must stay below $2^{30}$ signatures per key.

\paragraph{One-time layer.} Each \wc{} key only ever signs its own FORS public key, which is deterministic,
so reusing a leaf never reveals a second message to that key. Because all valid digests have the same digit
sum, no valid digest dominates another. A forgery therefore needs exactly the same digest, i.e.\ a second
preimage of the 128-bit value $\mathbf{T}_C$~\cite{sphincsc}.

\section{Signing modes and costs}
\label{sec:costs}

\paragraph{Cost model.} Costs count SHA-256 compression calls in the FIPS~205 SHA2 layout, with the
$\mathsf{PK.seed}$ block precomputed:
\begin{itemize}
  \item $\mathbf{PRF}$, $\mathbf{F}$, $\mathbf{H}$ and $\mathbf{T}_C$ cost 1;
        $\mathbf{T}_\ell$ costs $\lceil (22 + 16\ell + 9)/64 \rceil$.
  \item $\mathbf{H}_{msg}$ costs 5 for a short or pre-hashed message.
  \item One message attempt ($\mathbf{PRF}_{msg}$ and $\mathbf{H}_{msg}$) costs 7.
  \item Generating one leaf (64 PRF, 192 chain steps, $\mathbf{T}_{64}$, one tree node) costs 274.
\end{itemize}

\paragraph{Full keys.} Key generation builds all $2^{26}$ leaves: $18.4\cdot10^9$ compressions.
The signer keeps a 1~MiB cache ($2^{16}$ nodes) holding level~10 of the tree. Each signature then:
\begin{itemize}
  \item rebuilds the $2^{10}$ leaves under its cached node ($2.8\cdot10^5$);
  \item hashes the cached level up to the root for the upper authentication path ($6.6\cdot10^4$);
  \item signs with FORS ($7.4\cdot10^4$) and \wc{} ($1.0\cdot10^3$).
\end{itemize}
The total is $4.2\cdot10^5$ compressions.

\paragraph{Pruned keys.} Following hypertree pruning~\cite{conduition}, a weak signer computes only $K$
real leaves. It replaces every other leaf, or whole subtree of pruned leaves, by a pseudorandom surrogate
node derived from its secret seed. Key generation then costs about $274K$ compressions, and the whole kept
tree fits in the cache.
\begin{itemize}
  \item \textbf{Signing.} To sign, the signer grinds the randomizer $R$ until the leaf index lands on a real
        leaf. That takes $2^{26}/K$ attempts on average. It then signs with that leaf's FORS and \wc{} keys.
  \item \textbf{Security.} The analysis above holds with $2^h$ replaced by $K$. A key keeping $K$ leaves
        therefore signs $2^{30.11}\cdot K/2^{26}$ messages at 128-bit security:
\end{itemize}
\begin{center}
\begin{tabular}{lcccc}
\toprule
Pruned lifetime & real leaves $K$ & key generation & signing & of which grinding \\
\midrule
$2^{16}$ & $2^{11.9} \approx 3\,800$ & $1.04\cdot10^6$ & $2.0\cdot10^5$ & $1.2\cdot10^5$ \\
$2^{18}$ & $2^{13.9}$ & $4.2\cdot10^6$ & $1.1\cdot10^5$ & $3.1\cdot10^4$ \\
$2^{20}$ & $2^{15.9}$ & $1.7\cdot10^7$ & $8.3\cdot10^4$ & $7.7\cdot10^3$ \\
\bottomrule
\end{tabular}
\end{center}
Pruned signers must follow three rules.
\begin{itemize}
  \item \textbf{Separate seeds.} A pruned key must use a secret seed that is domain-separated from full
        keys, e.g.\ $\mathsf{SK.seed}' = \mathbf{PRF}(\mathsf{SK.seed}, \text{``pruned''})$. Otherwise a pruned and
        a full signer sharing a seed could make one \wc{} key sign two different FORS keys.
  \item \textbf{Fixed pattern.} The set of real leaves is chosen pseudorandomly from the seed at key
        generation and never changes.
  \item \textbf{Visibility.} Signatures remain valid, but after many signatures an observer can notice
        that only some leaves are used.
\end{itemize}

\section{Rationale}

\begin{itemize}
  \item \textbf{Why $h = 26$.} For a pruned key,
        (key generation) $\times$ (grinding) $\approx 2^{h} \cdot 274 \cdot 7$, independently of FORS. Targeting
        about $10^6$ for key generation and about $2^{17}$ for signing gives $h \approx 26$.
  \item \textbf{Why $a = 10$, $k = 24$.} The pruned signing budget leaves less than $10^5$ compressions
        for FORS, which forces small trees ($3k\cdot2^{a}$). The smallest $k$ that reaches $2^{30}$ at $h = 26$
        is then $k = 24$ for $a = 10$. That takes 4\,224 of the 5\,684 signature bytes.
  \item \textbf{Why a single tree.} It keeps the \slh{}-128-24 structure and the cheapest verification.
        The cost is heavy key generation for full keys, which pruning avoids for weak signers.
  \item \textbf{Why \wc{}.} Verification becomes constant: 366 hash calls for every signature, instead of
        305--494 for WOTS$^+$ with $w = 4$. That matters when verifying inside a proof system.
        It also saves \bytes{60}.
\end{itemize}

\begin{center}
\begin{tabular}{lcc}
\toprule
 & \slh{}-128-24 & \lean{} \\
\midrule
Signature & \bytes{3\,856} & \bytes{5\,684} \\
Lifetime at 128 bits & $2^{24.14}$ & $2^{30.11}$ \\
Verification (hash calls) & 187--376 & 366 \\
Full key generation & $1.2\cdot10^9$ & $1.8\cdot10^{10}$ \\
Pruned key ($2^{16}$ signatures): key generation / signing & $4.3\cdot10^6$ / $3.0\cdot10^8$ & $1.0\cdot10^6$ / $2.0\cdot10^5$ \\
\bottomrule
\end{tabular}
\end{center}

\section{Deviations from FIPS~205 and open points}

\begin{itemize}
  \item \textbf{Verifier changes.} New parameters (FIPS~205 does not allow them), and \wc{} in place of
        WOTS$^+$, which needs the counter field, the new address type and the digit-sum check. The
        verifier is otherwise standard.
  \item \textbf{Proof.} Security relies on the generic-attack estimate above. There is no tight proof in
        the quantum random-oracle model; Kudinov and Nick~\cite{kudinov-nick} estimate a loss of about
        10 bits.
  \item \textbf{Grinding assumption.} The grinding costs assume short or pre-hashed messages, since
        $\mathbf{H}_{msg}$ must be recomputed for each attempt.
  \item \textbf{Counter size.} A 16-bit \wc{} counter would suffice (failure probability about $e^{-79}$)
        and would save \bytes{2}.
\end{itemize}

\begin{thebibliography}{9}
\bibitem{fips205} NIST. \emph{Stateless Hash-Based Digital Signature Standard}. FIPS~205, 2024.
\bibitem{sp800-230} Q.~Dang, D.~Moody. \emph{Additional SLH-DSA Parameter Sets for Limited-Signature Use
  Cases}. NIST SP~800-230 (initial public draft), 2026.
\bibitem{sphincsc} M.~Kudinov, A.~H\"ulsing, E.~Ronen, E.~Yogev. \emph{SPHINCS$^+$C: Compressing SPHINCS$^+$
  With (Almost) No Cost}. IEEE S\&P, 2023.
\bibitem{sphincsplus} D.~J.~Bernstein et al. \emph{SPHINCS$^+$ submission to the NIST PQC project}, 2020.
\bibitem{fenner} C.~Fenner. \emph{slh-dsa-rls}: parameter search for SLH-DSA.
  \url{https://github.com/chrisfenner/slh-dsa-rls}, 2025.
\bibitem{kudinov-nick} M.~Kudinov, J.~Nick. \emph{Hash-based Signature Schemes for Bitcoin}. Cryptology ePrint
  Archive 2025/2203.
\bibitem{conduition} conduition. \emph{Pruning Hypertrees for the Lax and Lazy}.
  \url{https://conduition.io/cryptography/hypertree-pruning/}, 2026.
\end{thebibliography}

\end{document}
